\section*{Chapitre \ref{ch:g12:adaptive-immunity} --- L'immunité adaptative et la vaccination}

\begin{solution}{exo:g12:adaptive-immunity:1}
Un \omterm{def:g12:adaptive-immunity:antigen}{antigène} est toute molécule que le système adaptatif sait
reconnaître~; un \omterm{def:g12:adaptive-immunity:antigen}{lymphocyte} est un globule blanc portant des
récepteurs pour un seul \omterm{def:g12:adaptive-immunity:antigen}{antigène}. Sa spécificité vient de son
récepteur, assemblé au hasard pendant la maturation et identique sur
tous ses exemplaires.
\end{solution}

\begin{solution}{exo:g12:adaptive-immunity:2}
L'\omterm{def:g12:adaptive-immunity:antigen}{antigène} lie les quelques \omterm{def:g12:adaptive-immunity:antigen}{lymphocytes} dont le récepteur lui
convient~; ces \omterm{def:g10:cells-common-unit:cell}{cellules} se multiplient en un clone au fil de quelques
jours~; le clone se différencie en \omterm{def:g10:cells-common-unit:cell}{cellules} effectrices qui
combattent~; une minorité devient des \omterm{prop:g12:adaptive-immunity:selection}{cellules mémoire} qui persistent.
\end{solution}

\begin{solution}{exo:g12:adaptive-immunity:3}
Une \omterm{def:g11:gene-expression:protein}{protéine} en forme d'Y sécrétée par les plasmocytes, copie soluble
du récepteur du \omterm{def:g12:adaptive-immunity:antigen}{lymphocyte} B~; chaque bras lie l'\omterm{def:g12:adaptive-immunity:antigen}{antigène}. Sur une
bactérie, il forme un enrobage que les \omterm{def:g12:innate-immunity:phagocytosis}{phagocytes} saisissent, et il
ponte les bactéries en complexes qui sont englouties.
\end{solution}

\begin{solution}{exo:g12:adaptive-immunity:4}
Les \omterm{prop:g12:adaptive-immunity:tcells}{lymphocytes T auxiliaires}, activés par les \omterm{prop:g12:innate-immunity:handover}{cellules dendritiques},
sécrètent les cytokines qui autorisent et amplifient les \omterm{def:g12:adaptive-immunity:antigen}{lymphocytes} B
et les autres \omterm{def:g12:adaptive-immunity:antigen}{lymphocytes} T. Les \omterm{def:g12:adaptive-immunity:antigen}{lymphocytes} T cytotoxiques tuent les
\omterm{def:g10:cells-common-unit:cell}{cellules} qui exposent des fragments d'un virus ou d'une \omterm{def:g11:gene-expression:protein}{protéine}
tumorale.
\end{solution}

\begin{solution}{exo:g12:adaptive-immunity:5}
Une forme inoffensive de l'\omterm{def:g12:adaptive-immunity:antigen}{antigène} de l'agent pathogène (atténué,
tué, une \omterm{def:g11:gene-expression:protein}{protéine}, une toxine inactivée, ou ses instructions), avec un
adjuvant. Il provoque une réaction primaire et laisse des \omterm{prop:g12:adaptive-immunity:selection}{cellules
mémoire}, si bien que l'agent véritable rencontre une réaction
secondaire.
\end{solution}

\begin{solution}{exo:g12:adaptive-immunity:6}
Primaire~: une semaine de délai, un maximum bas (environ le dixième de
celui de la secondaire) à deux ou trois semaines, une décroissance en
un mois. Secondaire~: deux ou trois jours de délai, un maximum dix
fois plus haut en dix jours, qui persiste des mois.
\end{solution}

\begin{solution}{exo:g12:adaptive-immunity:7}
La mémoire est spécifique~: les \omterm{prop:g12:adaptive-immunity:selection}{cellules mémoire} laissées par la
première dose de A ne répondent qu'à A~; B, rencontré pour la première
fois, obtient une réaction primaire.
\end{solution}

\begin{solution}{exo:g12:adaptive-immunity:8}
Les \omterm{prop:g12:adaptive-immunity:antibodies}{anticorps} sont des \omterm{def:g10:chemistry-of-life:families}{protéines} des liquides et ne peuvent pas entrer
dans les \omterm{def:g10:cells-common-unit:cell}{cellules}. Une \omterm{def:g10:cells-common-unit:cell}{cellule} infectée expose à sa surface des
fragments du virus, et les \omterm{def:g12:adaptive-immunity:antigen}{lymphocytes} T cytotoxiques les
reconnaissent et tuent la \omterm{def:g10:cells-common-unit:cell}{cellule}.
\end{solution}

\begin{solution}{exo:g12:adaptive-immunity:9}
$1 - 1/4 = 75\%$~; $1 - 1/18 \approx 94\%$.
\end{solution}

\begin{solution}{exo:g12:adaptive-immunity:10}
Au-dessus du seuil, chaque cas en infecte moins d'un autre, si bien
que le virus ne peut pas circuler et n'atteint que rarement un
nourrisson~; les adultes immunisés sont le mur qui les entoure.
\end{solution}

\begin{solution}{exo:g12:adaptive-immunity:11}
Environ 750, 520 et 180 par microlitre~; le \omterm{prop:g12:adaptive-immunity:hiv}{sida} commence vers
l'année 9, quand le nombre passe sous 200.
\end{solution}

\begin{solution}{exo:g12:adaptive-immunity:12}
Les cytokines des \omterm{prop:g12:adaptive-immunity:tcells}{lymphocytes T auxiliaires} sont ce qui autorise les
\omterm{def:g12:adaptive-immunity:antigen}{lymphocytes} B à devenir plasmocytes et les \omterm{def:g12:adaptive-immunity:antigen}{lymphocytes} T cytotoxiques
à se multiplier~; les deux bras dépendent du même coordinateur, et le
supprimer les réduit tous deux au silence.
\end{solution}

\begin{solution}{exo:g12:adaptive-immunity:13}
Les \omterm{prop:g12:adaptive-immunity:antibodies}{anticorps} ont été sélectionnés contre le virus tel qu'il était~;
le virus mute ses \omterm{def:g10:chemistry-of-life:families}{protéines} de surface plus vite que de nouveaux
clones ne peuvent être sélectionnés, et chaque réaction est donc
distancée. Et le virus se cache et se multiplie dans les \omterm{def:g10:cells-common-unit:cell}{cellules}
auxiliaires mêmes qui coordonneraient la réaction.
\end{solution}

\begin{solution}{exo:g12:adaptive-immunity:14}
Les \omterm{prop:g12:adaptive-immunity:selection}{cellules mémoire} reconnaissent les \omterm{def:g10:chemistry-of-life:families}{protéines} de surface que le
\omterm{def:g12:adaptive-immunity:vaccine}{vaccin} a présentées. Celles de la grippe mutent chaque année en formes
que l'ancienne mémoire ne lie pas~; celles de la rougeole ne changent
presque pas, si bien que la mémoire de 1970 convient encore au virus
d'aujourd'hui.
\end{solution}

\begin{solution}{exo:g12:adaptive-immunity:15}
Ce qui varie~: les récepteurs des \omterm{def:g12:adaptive-immunity:antigen}{lymphocytes}, faits au hasard. Ce qui
sélectionne~: l'\omterm{def:g12:adaptive-immunity:antigen}{antigène}, qui ne laisse se multiplier que les clones
adaptés. Ce qui est hérité~: le récepteur, par les descendantes du
clone, \omterm{prop:g12:adaptive-immunity:selection}{cellules mémoire} comprises. Échelle de temps~: des jours pour
une réaction, une vie pour la mémoire, contre des générations pour une
\omterm{def:g12:selection-drift-speciation:population}{population}. L'analogie s'arrête à la transmission~: les clones
sélectionnés meurent avec le corps et ne passent pas à la génération
suivante.
\end{solution}

\begin{solution}{pb:g12:adaptive-immunity:1}
\textbf{1.} Rien pendant environ une semaine, un maximum modeste à
deux ou trois semaines, une décroissance au fil des semaines
suivantes.

\textbf{2.} Des \omterm{prop:g12:adaptive-immunity:antibodies}{anticorps} en deux ou trois jours, un maximum dix fois
plus haut, qui dure des mois~: les \omterm{prop:g12:adaptive-immunity:selection}{cellules mémoire} laissées par la
première dose forment déjà un clone nombreux, de sorte que la
multiplication part de milliers de \omterm{def:g10:cells-common-unit:cell}{cellules} et non d'une poignée.

\textbf{3.} Un virus vivant se multiplie brièvement dans les \omterm{def:g10:cells-common-unit:cell}{cellules},
si bien que les \omterm{def:g10:cells-common-unit:cell}{cellules} infectées exposent ses fragments et que des
\omterm{def:g12:adaptive-immunity:antigen}{lymphocytes} T cytotoxiques sont sélectionnés en plus des \omterm{prop:g12:adaptive-immunity:antibodies}{anticorps}~;
un virus tué entraîne surtout le bras des \omterm{prop:g12:adaptive-immunity:antibodies}{anticorps}.

\textbf{4.} Les \omterm{prop:g12:adaptive-immunity:antibodies}{anticorps} déjà présents bloquent l'essentiel du
virus~; les \omterm{prop:g12:adaptive-immunity:selection}{cellules mémoire} B et T se multiplient en deux jours~; les
\omterm{def:g10:cells-common-unit:cell}{cellules} infectées sont tuées avant que le virus se répande.
L'infection est arrêtée avant les symptômes.

\textbf{5.} Trois jours de multiplication silencieuse, puis fièvre,
éruption et maladie tandis que la réaction primaire se construit sur
une semaine~; guérison vers le dixième jour, avec des \omterm{prop:g12:adaptive-immunity:selection}{cellules mémoire}
et une immunité durable --- s'il survit aux complications.

\textbf{6.} $1 - 1/15 \approx 93\%$.

\textbf{7.} $15 \times 0.10 = 1.5$~: chaque cas en infecte plus d'un
autre, la maladie peut donc se propager. La ville n'est pas protégée.

\textbf{8.} $10\% \times \num{50000} = 5000$ non immunisés, dont 1500
sont les nourrissons.

\textbf{9.} 2500 personnes réceptives~; $15 \times 0.05 = 0.75$~:
moins d'un, la maladie s'éteint.

\textbf{10.} Les nourrissons ne peuvent pas être vaccinés, de sorte
que leur seule protection est que les adultes qui les entourent, en
étant immunisés, ne laissent au virus aucun chemin pour les atteindre.

\textbf{11.} Troisième génération~: $1.5^3 \approx 3.4$ cas~; sixième
génération~: $1.5^6 \approx 11$.

\textbf{12.} $1 + 1.5 + 2.25 + 3.4 + 5.1 + 7.6 + 11.4 \approx 32$ cas.

\textbf{13.} Environ 30\% de 32~: quelque 10 nourrissons.

\textbf{14.} Chaque cas retire une personne réceptive et en ajoute une
immunisée, de sorte que la fraction réceptive baisse, et avec elle le
nombre de personnes que chaque cas infecte~; avec 5000 personnes
réceptives, l'épidémie ne ralentirait qu'après des centaines de cas,
au bout de plusieurs mois --- à moins que la vaccination n'intervienne.

\textbf{15.} À 95\%~: $0.75^3 \approx 0.4$ et $0.75^6 \approx 0.18$
cas~; total d'environ $1 + 0.75 + 0.56 + \dots \approx 4$ cas. Dix
fois moins de cas, et l'épidémie s'éteint d'elle-même.

\textbf{16.} 800 en 8 ans~: 100 par microlitre et par an, soit environ
8 par mois.

\textbf{17.} Les \omterm{def:g12:adaptive-immunity:antigen}{lymphocytes} B ont besoin des cytokines des \omterm{def:g10:cells-common-unit:cell}{cellules}
auxiliaires pour devenir plasmocytes, et les \omterm{def:g12:adaptive-immunity:antigen}{lymphocytes} T
cytotoxiques en ont besoin pour se multiplier~: avec trop peu
d'auxiliaires, aucun des deux bras n'est autorisé, bien que leurs
\omterm{def:g10:cells-common-unit:cell}{cellules} existent.

\textbf{18.} Oui~: les \omterm{prop:g12:adaptive-immunity:selection}{cellules mémoire} ont encore besoin des
cytokines auxiliaires pour se déployer en réaction secondaire~; sans
coordination, la mémoire ne peut pas servir, et la rougeole,
inoffensive pour un adulte sain, peut être mortelle.

\textbf{19.} La \omterm{def:g12:selection-drift-speciation:population}{population} auxiliaire repousse à partir des \omterm{def:g10:cells-common-unit:cell}{cellules}
survivantes et de la moelle, et les \omterm{prop:g12:adaptive-immunity:selection}{cellules mémoire} des autres bras
n'ont jamais été détruites --- le répertoire est intact dès que la
coordination revient. Le virus persiste dans des \omterm{def:g10:cells-common-unit:cell}{cellules} cachées, si
bien qu'arrêter les médicaments le laisse repartir.

\textbf{20.} 93\% de couverture~; environ 32 cas à partir d'un seul
voyageur à 90\%~; le \omterm{prop:g12:adaptive-immunity:tcells}{lymphocyte T auxiliaire}.
\end{solution}
